\subsection*{Consistency}
\label{sec:axiom-systems-consistency}

\begin{tcolorbox}[colback=gray!6, colframe=gray!40, arc=2pt,
  left=8pt, right=8pt, top=6pt, bottom=6pt,
  title={\small\textbf{Consistency Presentation: A Simple Example}},
  fonttitle=\small\bfseries]
Consider two informal axiom systems.

\medskip
\noindent\textbf{Example A.}
\begin{itemize}
  \item All objects are red.
  \item Some object is not red.
\end{itemize}
These assertions cannot simultaneously hold. The first assertion says every
object has the property; the second says at least one object lacks it.

\medskip
\noindent\textbf{Example B.}
\begin{itemize}
  \item All objects are red.
  \item Object $a$ is red.
\end{itemize}
These assertions can hold simultaneously. The second assertion agrees with
what the first assertion already requires.

\medskip
These examples are only for intuition. A consistent axiom system is one whose
assertions fit together without contradiction. An inconsistent axiom system is
one whose assertions cannot be held together coherently.
\end{tcolorbox}

\begin{remark*}[Structural remarks]
\textbf{Why contradictions matter.}
Contradictory systems are unusable because their assertions do not fit
together. If a system both requires and forbids the same situation, then it no
longer gives a coherent basis for reasoning.

\textbf{Models as evidence.}
If a model exists, the axioms hold together coherently in at least one
interpretation. A model is therefore evidence that the axioms are mutually
compatible.

\textbf{Absolute vs relative consistency.}
Most consistency results in mathematics are relative consistency results. They
show that one theory is consistent provided some other background theory is
consistent. This section records the idea only; the technical development
belongs elsewhere.

\textbf{Forward roadmap.}
The next section uses models and consistency to introduce independence:
\[
\begin{array}{c}
\text{Consistency}\\
\Downarrow\\
\text{Independence}
\end{array}
\]
\end{remark*}

\begin{definition}[Contradiction]
\label{def:axiom-systems-contradiction}
A \emph{contradiction} is a statement that cannot be true under a given
interpretation.
\end{definition}

\begin{remark*}[Interpretation]
A contradiction marks a breakdown in compatibility. It identifies an assertion
that cannot be made true in the interpretation under consideration.
\end{remark*}

\begin{dependencies}
\begin{itemize}
  \item \hyperref[def:axiom-systems-satisfaction]{Satisfaction}
\end{itemize}
\end{dependencies}

\begin{definition}[Consistent Axiom System]
\label{def:axiom-systems-consistent-axiom-system}
A \emph{consistent axiom system} is an axiom system that does not lead to
contradiction.
\end{definition}

\begin{remark*}[Interpretation]
Consistency says that the adopted assertions fit together. A consistent axiom
system avoids forcing an impossible situation.
\end{remark*}

\begin{dependencies}
\begin{itemize}
  \item \hyperref[def:axiom-system]{Axiom System}
  \item \hyperref[def:axiom-systems-contradiction]{Contradiction}
\end{itemize}
\end{dependencies}

\begin{definition}[Inconsistent Axiom System]
\label{def:axiom-systems-inconsistent-axiom-system}
An \emph{inconsistent axiom system} is an axiom system containing or producing
contradictory assertions.
\end{definition}

\begin{remark*}[Interpretation]
Inconsistency means that the adopted assertions cannot all be maintained
coherently. The system asks for incompatible requirements at once.
\end{remark*}

\begin{dependencies}
\begin{itemize}
  \item \hyperref[def:axiom-system]{Axiom System}
  \item \hyperref[def:axiom-systems-contradiction]{Contradiction}
\end{itemize}
\end{dependencies}

\begin{definition}[Model Existence Criterion (Informal)]
\label{def:axiom-systems-model-existence-criterion}
Informally, if an axiom system has a model, then the axioms are mutually
compatible.
\end{definition}

\begin{remark*}[Interpretation]
This is an informal criterion in this section, not a formal theorem. The idea
is that a model provides a setting in which all the axioms hold together, so
the axioms have at least one coherent realization.
\end{remark*}

\begin{dependencies}
\begin{itemize}
  \item \hyperref[def:axiom-systems-model]{Model}
\end{itemize}
\end{dependencies}

\begin{definition}[Relative Consistency]
\label{def:axiom-systems-relative-consistency}
Let $T_0$ and $T_1$ be theories. The theory $T_1$ is
\emph{consistent relative to $T_0$} if
\[
  \operatorname{Con}(T_0)\implies \operatorname{Con}(T_1).
\]
Equivalently, every contradiction in $T_1$ would imply a contradiction in
$T_0$.
\end{definition}

\begin{remark*}[Standard quantified statement]
For theories $T_0$ and $T_1$,
\[
  \operatorname{RelCon}(T_1,T_0)
  \quad\Longleftrightarrow\quad
  \bigl(\neg\exists d_0\,(d_0:T_0\vdash\bot)\bigr)
  \implies
  \bigl(\neg\exists d_1\,(d_1:T_1\vdash\bot)\bigr).
\]
Equivalently, in contrapositive form,
\[
  \exists d_1\,(d_1:T_1\vdash\bot)
  \implies
  \exists d_0\,(d_0:T_0\vdash\bot).
\]
\end{remark*}

\begin{remark*}[Interpretation]
Relative consistency does not establish consistency from nothing. It says that
the target theory introduces no new contradiction unless the base theory was
already inconsistent. In practice, a relative consistency proof often gives a
translation that turns any contradiction derived in $T_1$ into a contradiction
derived in $T_0$.
\end{remark*}

\begin{dependencies}
\begin{itemize}
  \item \hyperref[def:axiom-systems-consistent-axiom-system]{Consistent Axiom System}
\end{itemize}
\end{dependencies}
